\section{Related Work}
\label{sec:related}

\paragraph{Automated review, and grading by prompting.}
Using LLMs to assist peer review has moved from speculation to deployment: a randomized
study at ICLR~2025 improved review informativeness with an LLM feedback agent
\citep{thakkar2026feedback}, and AAAI~2026 ran an AI-review pilot at scale
\citep{biswas2026aaai}. The anchor empirical study found GPT-4 review comments overlap
human comments about as much as two humans overlap \citep{liang2024llm}; multi-agent and
pipeline systems raise specificity \citep{darcy2024marg,lu2024aiscientist}; recent
benchmarks measure LLM reviewer ability and its divergence from humans
\citep{li2026llmreviewer,reviewertoo2025}; and critical work documents brittleness,
hidden-prompt attacks, and inflated scores
\citep{du2024llmsassist,zhou2024reliable,lin2025hidden,zhou2025positivereview,russolatona2024lottery,pangram2026iclr},
prompting funding agencies and journals to restrict undisclosed AI review
\citep{science2025funding}. Almost all of this evaluates generated review \emph{text}.
Separately, \emph{grading by prompting} is an established paradigm: strong LLMs score
open-ended text against rubrics with near-human agreement (LLM-as-a-judge;
\citealp{zheng2023judging,liu2023geval}), and prompt-only scoring has been validated
against human grades in automated essay scoring \citep{mizumoto2023exploring}. We build
on that paradigm rather than claim it.

\paragraph{Noisy ground truth, and how to validate against it.}
Any claim that a signal tracks ``quality'' must be calibrated against the unreliability
of the human ground truth. The NeurIPS consistency experiments are canonical: roughly
half of accepted papers would change under a re-run, and review succeeds at rejecting
weak papers but struggles to rank strong ones
\citep{cortes2021inconsistency,beygelzimer2023arbitrary}; bias and its measurement are
well documented \citep{tomkins2017reviewerbias,stelmakh2019biases,shah2022challenges}.
For validating a score against noisy, ordinal outcomes the field uses rank correlation,
threshold-free AUROC, and calibration assessment \citep{spearman1904,guo2017calibration},
on corpora such as PeerRead and NLPeer \citep{kang2018peerread,dycke2023nlpeer};
\citet{checco2021aiassisted} validated a learned score against decisions and ratings,
though with a trained classifier rather than an LLM.

\paragraph{Positioning.}
The older paradigm for scholarly assessment \emph{trains} on reviews and decisions
(acceptance prediction; \citealp{kang2018peerread}) or \emph{fine-tunes} review
generators \citep{idahl2025openreviewer}. Prompt-only treatment of manuscripts has so
far targeted feedback \emph{text} \citep{liang2024llm} or has been validated only against
weak proxies, such as an author's own ratings of his own papers
\citep{thelwall2024chatgpt}. To our knowledge this is the first study to validate a
training-free, prompt-only manuscript quality \emph{score} against the public \emph{reject}
class and reviewer ratings, and to decompose how much of that validity comes from the
underlying model versus the engineered pipeline. The individual ingredients are not each
new --- prompting, pre-registration, and outcome validation all have precedent --- but
their combination on the observable reject class, with the model-versus-pipeline
reliability finding, is. The novelty is the validity result, not the prompting.
